%% kapa: exam=1 round=32 year=2021 period=1 form=A subject=real_estate subject_ko=부동산학원론 seq=5 number=85
%% source: Q-Net 감정평가사 기출문제 (cst00309) artlSeq=5208397 file=extracted/1008/5208397/A형/2021년도 제32회 감정평가사 1차 1교시 A형.pdf
%% conversion: kordoc 4.24.51 -> tools/gen1.py (auto); review=proofread-sonnet; flags=none
\begin{question}{85}
수요함수와 공급함수가 각각 A부동산시장에서는 $Q_d=200-P$, $Q_s=10+\dfrac{1}{2}P$이고 B부동산시장에서는 $Q_d=400-\dfrac{1}{2}P$, $Q_s=50+2P$이다. 거미집이론(Cob-web theory)에 의한 A시장과 B시장의 모형 형태의 연결이 옳은 것은? (단, x축은 수량, y축은 가격, 각각의 시장에 대한 $P$는 가격, $Q_d$는 수요량, $Q_s$는 공급량이며, 가격변화에 수요는 즉각 반응하지만 공급은 시간적인 차이를 두고 반응함, 다른 조건은 동일함)
\begin{choices}
  \item A: 발산형, B: 수렴형
  \item A: 발산형, B: 순환형
  \item A: 순환형, B: 발산형
  \item A: 수렴형, B: 발산형
  \item A: 수렴형, B: 순환형
\end{choices}
\end{question}
\answer{4}
